\documentclass[11pt]{amsart}
\usepackage[margin=1.15in]{geometry}
\usepackage{amsmath,amssymb,amsthm,mathtools,enumitem}
\usepackage{microtype}
\usepackage{xcolor}
\usepackage[colorlinks=true,linkcolor=blue!50!black,citecolor=blue!50!black]{hyperref}

\newtheorem{theorem}{Theorem}[section]
\newtheorem{lemma}[theorem]{Lemma}
\newtheorem{proposition}[theorem]{Proposition}
\newtheorem{conjecture}[theorem]{Conjecture}
\theoremstyle{definition}
\newtheorem{definition}[theorem]{Definition}
\newtheorem{fact}[theorem]{Fact}
\theoremstyle{remark}
\newtheorem{remark}[theorem]{Remark}

\newcommand{\T}{\mathbb{T}}
\newcommand{\Pl}{\mathbb{P}}
\newcommand{\Iso}{\operatorname{Iso}}
\newcommand{\arc}[2]{(#1,#2)}
\newcommand{\cl}{\operatorname{cl}}
\newcommand{\struct}{\operatorname{struct}}
\newcommand{\cost}{\operatorname{cost}}
\newcommand{\Prig}{\mathcal{P}_{\mathrm{rig}}}
\newcommand{\Surv}{\mathcal{P}}
\newcommand{\up}{H^{\uparrow}}
\newcommand{\down}{H^{\downarrow}}

\title[Sequential realizability of a flat vertex]{Sequential realizability of a flat vertex:\\
a rigid layered model and one six-crease configuration}
\date{}

\begin{document}
\begin{abstract}
We fix a combinatorial model of folding a single flat vertex by a sequence of steps, each of
which folds along one line of the current folded image and may create and erase creases. We
add a layered refinement in which each step is a rigid half-turn of a block of layers. For the
vertex $A=\{0,1,3,4,5,7\}\cdot\pi/4$ we show that a natural history with cancellation, namely
\emph{fold both diagonals, fold the full cross, unfold the line $\{\pi/2,3\pi/2\}$}, is valid in the base
model and passes the collision test (C), but admits no rigid realization. We introduce the
\emph{surviving partition} of a history, which groups the final creases by the step that last
created them and forgets the order of the steps. We conjecture that no rigidly realizable
history of $A$ has the surviving partition $\{\{0,4\},\{1,5\},\{3,7\}\}$ (three straight lines), nor
indeed any surviving partition into straight lines through the vertex.
\end{abstract}
\maketitle


\section{Base model}

\subsection*{Conventions}
$\T := \mathbb{R}/2\pi\mathbb{Z}$ (directions), ordered via representatives in $[0,2\pi)$;
$\Pl := \mathbb{R}/\pi\mathbb{Z}$ (lines) with projection $q\colon \T\to\Pl$.
$\arc{a}{b}$ denotes the open counterclockwise arc from $a$ to $b$, and $\cl$ its closure.
For $c\in\T$ let $\sigma_c(x) := 2c-x$. Since $2(c+\pi)=2c$ in $\T$, $\sigma_c$ depends only on
$q(c)$, so $\sigma_\ell$ is defined for $\ell\in\Pl$.
$\Iso(\T) := \{x\mapsto \varepsilon x+t : \varepsilon=\pm1,\ t\in\T\}$.

\begin{fact}\label{fact:conj}
For $g\in\Iso(\T)$ and $x\in\T$: $g\circ\sigma_x\circ g^{-1}=\sigma_{g(x)}$.
\end{fact}
\begin{proof}
For $g(y)=\varepsilon y+t$: $g\sigma_x g^{-1}(y)=\varepsilon\bigl(2x-\varepsilon(y-t)\bigr)+t
=2(\varepsilon x+t)-y$.
\end{proof}

\begin{definition}[Faces, folding maps, flatness]\label{def:flat}
Let $S\subset\T$ be finite, $S=\{s_1<\dots<s_m\}$. If $m\ge1$ the \emph{faces} of $S$ are
$F_j:=\arc{s_j}{s_{j+1}}$ (indices mod $m$); if $S=\emptyset$ the only face is $\T$.
The \emph{folding maps} are $\varphi_1:=\mathrm{id}$ and
$\varphi_{j+1}:=\varphi_j\circ\sigma_{s_{j+1}}$ for $1\le j<m$.
$S$ is \emph{flat} if $S=\emptyset$ or $\varphi_m\circ\sigma_{s_1}=\mathrm{id}$.
\end{definition}

\begin{fact}\label{fact:Phi}
Let $S\ne\emptyset$ be flat. Then $m$ is even and every face has length $\le\pi$ (Kawasaki).
The map $\Phi_S\colon\T\to\T$, $\Phi_S|_{\cl F_j}:=\varphi_j$, is well defined and continuous,
since $\varphi_{j+1}(s_{j+1})=\varphi_j\sigma_{s_{j+1}}(s_{j+1})=\varphi_j(s_{j+1})$ and,
at $s_1$, $\varphi_m(s_1)=\varphi_m\sigma_{s_1}(s_1)=s_1=\varphi_1(s_1)$.
Set $\Phi_\emptyset:=\mathrm{id}$.
\end{fact}

\begin{definition}\label{def:Pi}
For flat $S$ and $\ell\in\Pl$ let $\Pi_S(\ell):=\{x\in\T : q(\Phi_S(x))=\ell\}$.
It is finite: $\Phi_S$ is an isometry on each closed face, which has length $\le\pi$, so each
closed face contains at most two points of $\Pi_S(\ell)$; and $\Pi_\emptyset(\ell)=q^{-1}(\ell)$.
\end{definition}

\begin{definition}[Step]\label{def:step}
A \emph{step} from a flat $S$ is a nonempty $D\subset\T$ with $|D|$ even and
$D\subseteq\Pi_S(\ell)$ for some $\ell\in\Pl$; this $\ell$ is unique and denoted $\ell(D)$.
Put $N(D):=D\setminus S$ (created), $K(D):=D\cap S$ (erased). The result is $S\mathbin{\triangle}D$.
This is exactly the original definition: (P) is $|D|$ even, and (L) is $D\subseteq\Pi_S(\ell)$.
\end{definition}

\begin{definition}[Pieces and labelings]\label{def:label}
The \emph{pieces} of a step $D$ from $S$ are the faces of $S\cup D$. A \emph{labeling} is a map
$\mu\colon\{\text{pieces}\}\to\{0,1\}$ such that two consecutive pieces receive different labels
iff their common endpoint lies in $D$. As $|D|$ is even there are exactly two labelings, and
they are complementary. Label $1$ means \emph{moving}. The parent face of a piece $p$, i.e.\ the
face of $S$ containing $p$, is denoted $\pi(p)$.
\end{definition}

\begin{lemma}\label{lem:result}
Let $D$ be a step from flat $S$, $\ell=\ell(D)$, and $\mu$ a labeling. Define
$\Psi_\mu:=\Phi_S$ on closures of pieces with $\mu=0$ and $\Psi_\mu:=\sigma_\ell\circ\Phi_S$ on
closures of pieces with $\mu=1$. Then $S\mathbin{\triangle}D$ is flat, and
$\Phi_{S\triangle D}=g\circ\Psi_\mu$ for some $g\in\Iso(\T)$.
\end{lemma}
\begin{proof}
$\Psi_\mu$ is well defined and continuous because $\sigma_\ell$ fixes
$\Phi_S(D)\subseteq q^{-1}(\ell)$ pointwise. Compare the restrictions of $\Psi_\mu$ to the two
pieces meeting at a point $x\in S\cup D$. Below we treat the case where the piece before $x$ has
label $0$; the other case follows by composing both restrictions on the left with $\sigma_\ell$.
If $S=\emptyset$, read $\varphi_j=\mathrm{id}$.
\begin{itemize}[leftmargin=*]
\item If $x\in S\setminus D$, $x=s_{j+1}$, the labels agree and the restrictions are $\varphi_j$
and $\varphi_{j+1}=\varphi_j\sigma_x$.
\item If $x\in N(D)$ lies in a face $F_j$, then $q(\varphi_j(x))=\ell$, so by
Fact~\ref{fact:conj} $\varphi_j\sigma_x=\sigma_{\varphi_j(x)}\varphi_j=\sigma_\ell\varphi_j$.
So the restrictions are $\varphi_j$ and $\sigma_\ell\varphi_j=\varphi_j\sigma_x$, and $x$ becomes
a fold.
\item If $x\in K(D)$, $x=s_{j+1}$, then
the restrictions are $\varphi_j$ and
$\sigma_\ell\varphi_{j+1}=\sigma_\ell\varphi_j\sigma_x=\sigma_\ell\sigma_\ell\varphi_j=\varphi_j$,
using $\varphi_j\sigma_x=\sigma_\ell\varphi_j$ as before. So they coincide and the fold at $x$
disappears.
\end{itemize}
Thus $\Psi_\mu$ is continuous, piecewise in $\Iso(\T)$, and changes by $\sigma_x$ exactly at the
points $x$ of $S\mathbin{\triangle}D$. Let $g\in\Iso(\T)$ make $g\circ\Psi_\mu$ the identity on the first face of
$S\mathbin{\triangle}D$. Then the folding maps of $S\mathbin{\triangle}D$ are the restrictions of
$g\circ\Psi_\mu$, and the closing condition holds because $g\circ\Psi_\mu$ is continuous around
the whole circle. (If $S\mathbin{\triangle}D=\emptyset$, $\Psi_\mu$ is a single isometry and
$g=\Psi_\mu^{-1}$.)
\end{proof}

\begin{lemma}[Reversibility]\label{lem:rev}
If $D$ is a step from $S$, then $D$ is a step from $S\mathbin{\triangle}D$.
\end{lemma}
\begin{proof}
For $d\in D$, $\Psi_\mu(d)=\Phi_S(d)$ because $\sigma_\ell$ fixes $q^{-1}(\ell)$ pointwise. By
Lemma~\ref{lem:result}, $\Phi_{S\triangle D}(d)=g(\Phi_S(d))\in g\bigl(q^{-1}(\ell)\bigr)$, a single
line, for all $d\in D$.
\end{proof}

\begin{definition}[Histories]\label{def:hist}
A \emph{history} is a tuple $H=(D_1,\dots,D_n)$ with $S_0=\emptyset$,
$S_i=S_{i-1}\mathbin{\triangle}D_i$, and each $D_i$ a step from $S_{i-1}$. It is a history of $A$
if $S_n=A$. For $a\in A$ let $\tau_H(a):=\max\{i : a\in N(D_i)\}$; the set is nonempty since $a\notin S_0$. The \emph{structure}
$\struct(H)$ is the preorder $a\preceq b\iff\tau_H(a)\le\tau_H(b)$. The \emph{cost} is
$\cost(H):=\sum_i|K(D_i)|$. $R(A)$ is the set of structures of histories of $A$, and $R_0(A)$
the subset coming from histories of cost $0$.
\end{definition}

\begin{definition}[Surviving partition, straight blocks]\label{def:surv}
The \emph{surviving partition} of a history $H=(D_1,\dots,D_n)$ of $A$ is
\[
\Surv(H):=\bigl\{\tau_H^{-1}(i) : 1\le i\le n\bigr\}\setminus\{\emptyset\},
\]
a partition of $A$. Its blocks are called the \emph{surviving blocks} of $H$: the block
$\tau_H^{-1}(i)$ is what survives in $A$ of the creations made by step $i$. A set $B\subseteq\T$ is
\emph{straight} if $B+\pi=B$, i.e.\ $B$ is a union of lines through the vertex.
\end{definition}

\begin{remark}\label{rem:surv}
$\Surv(H)$ is the partition underlying $\struct(H)$, with the order of the classes forgotten.
It ignores scaffold creases, i.e.\ creases created and later erased, and it ignores earlier
creations of creases of $A$ that were erased and created again: only the last creation counts.
\end{remark}

\begin{definition}[Collision condition (C)]\label{def:C}
A step $D$ from $S$ satisfies (C) if the following holds. Whenever $F,F'$ are faces of $S$ with
$F\cap D=F'\cap D=\emptyset$, and $x\in F$, $x'\in F'$ satisfy
$\Phi_S(x)=\Phi_S(x')\in q^{-1}(\ell(D))$, then $\mu(F)=\mu(F')$. Such faces are pieces, and the
condition does not depend on which labeling is used.
\end{definition}

\section{Layered model}

\begin{definition}[Layered state]\label{def:layer}
Let $S$ be flat. For a face $F$ put $W_F:=\Phi_S(F)$, an open arc (or $\T$ if $S=\emptyset$). Let $E_S:=\Phi_S(S)$, a
finite set, and for $u\notin E_S$ let $L(u):=\{F : u\in W_F\}$.
A \emph{layer order} on $S$ is a family $(<_u)_{u\in\T\setminus E_S}$ of strict total orders on
$L(u)$ that is locally constant in $u$.

For $v\in E_S$, write $<_{v^+}$ and $<_{v^-}$ for the order at points of $\T\setminus E_S$
immediately counterclockwise, resp.\ clockwise, of $v$; likewise $L(v^\pm)$. Put
$\Theta(v):=L(v^-)\cap L(v^+)$, the faces \emph{passing through} $v$. The \emph{taco} of a ray
$s=s_{j+1}\in S$ is $\{F_j,F_{j+1}\}$. Since $\varphi_{j+1}(s+t)=\varphi_j(s-t)$, both faces of the
taco of $s$ lie in $L(v^{\epsilon})$ for exactly one sign $\epsilon=\epsilon(s)$, where
$v=\Phi_S(s)$. The pair $(S,<)$ is a \emph{layered state} if the following three conditions hold.
\begin{enumerate}[label=(O\arabic*)]
\item For every $v\in E_S$, the orders $<_{v^-}$ and $<_{v^+}$ agree on $\Theta(v)$.
\item If $s\ne s'$, $\Phi_S(s)=\Phi_S(s')=v$ and $\epsilon(s)=\epsilon(s')=\epsilon$, then the two
tacos, which are disjoint, do not interleave in $<_{v^\epsilon}$: it is not the case that
exactly one face of one taco lies strictly between the two faces of the other.
\item If $v=\Phi_S(s)$ and $c\in\Theta(v)$, then $c$ does not lie strictly between the two faces of
the taco of $s$ in $<_{v^{\epsilon(s)}}$.
\end{enumerate}
The \emph{initial layered state} is $\emptyset$ with its unique layer order.
\end{definition}

\begin{definition}[Rigid step]\label{def:rigid}
The data are a layered state $(S,<)$, a step $D$ with $\ell=\ell(D)$, a labeling $\mu$, and a
choice of one component $\up$ of $\T\setminus q^{-1}(\ell)$; the other component is $\down$.
Physically: a half-turn about $\ell$ that lifts the layers over $\up$ and lowers those over
$\down$.

\smallskip\noindent\textbf{(F)} Let $u\notin E_S\cup q^{-1}(\ell)$, and let $p,p'$ be pieces with
$\mu(p)=1$, $\mu(p')=0$ and $u\in\Phi_S(p)\cap\Phi_S(p')$. Then
$\pi(p')<_u\pi(p)$ if $u\in\up$, and $\pi(p)<_u\pi(p')$ if $u\in\down$.

\smallskip\noindent\textbf{New order.} Let $S'=S\mathbin{\triangle}D$ and let $\Psi=\Psi_\mu$ and
$g$ be as in Lemma~\ref{lem:result}. Each face $G$ of $S'$ is a union of pieces and points of
$K(D)$, and $\Psi$ is injective on $G$. Let $u$ avoid the finite set
$E_S\cup\sigma_\ell(E_S)\cup q^{-1}(\ell)$. Define $\prec_u$ on the pieces $p$ with
$u\in\Psi(p)$ as follows:
\begin{enumerate}[label=(\roman*)]
\item two staying pieces: $p\prec_u p'\iff\pi(p)<_u\pi(p')$;
\item two moving pieces: $p\prec_u p'\iff\pi(p')<_{\sigma_\ell(u)}\pi(p)$, i.e.\ the order is
reversed;
\item $p$ moving and $p'$ staying: $p'\prec_u p$ if $u\in\down$, and $p\prec_u p'$ if $u\in\up$.
\end{enumerate}
Set $G<'_{g(u)}G'$ iff $p_G\prec_u p_{G'}$, where $p_G$ is the unique piece of $G$ with
$u\in\Psi(p_G)$. The step is \emph{rigid} if (F) holds and $<'$ extends to a layer order making
$(S',<')$ a layered state.
\end{definition}

A history is \emph{rigidly realizable} if there are layered states $(S_i,<^i)$, starting from the
initial one, such that each transition $(S_{i-1},<^{i-1})\to(S_i,<^i)$ is a rigid step for $D_i$,
for some labeling and some choice of $\up$. Let
\[
\Prig(A):=\{\Surv(H) : H \text{ a rigidly realizable history of } A\}
\]
be the set of \emph{rigidly realizable surviving partitions} of $A$.

\begin{lemma}\label{lem:FC}
(F) implies (C). Consequently every step of a rigidly realizable history satisfies (C).
\end{lemma}
\begin{proof}
Let $F,F',x,x'$ be as in Definition~\ref{def:C} with $w:=\Phi_S(x)=\Phi_S(x')$, and suppose
$\mu(F)=1\ne0=\mu(F')$. Since $x,x'$ are interior points of faces, $F,F'\in\Theta(w)$ if
$w\in E_S$, and otherwise $F,F'\in L(u)$ for all $u$ near $w$ on both sides; in both cases their
relative order is the same on both sides of $w$ (by (O1), resp.\ local constancy). But (F)
requires $\pi(F')<_u\pi(F)$ for $u\in\up$ near $w$ and $\pi(F)<_u\pi(F')$ for $u\in\down$ near
$w$. Contradiction.
\end{proof}

\begin{lemma}[Simple folds]\label{lem:simple}
Let $D=\Pi_S(\ell)$ with $D\cap S=\emptyset$. Then the image of each piece lies in one component
of $\T\setminus q^{-1}(\ell)$, and each labeling moves exactly the pieces lying over one
component. Hence (F) holds vacuously, and for each of the four choices of labeling and $\up$ the order
$<'$ is determined by (i)--(iii).
\end{lemma}
\begin{proof}
No ray maps into $q^{-1}(\ell)$, since $\Pi_S(\ell)\cap S=D\cap S=\emptyset$, and all interior
points of faces mapping into $q^{-1}(\ell)$ lie in $D$. So no piece's image meets $q^{-1}(\ell)$
in its interior. Across a point of $D$ (a transversal crossing) both the side and the label
switch; across a ray of $S$ (a fold away from $\ell$) neither switches.
\end{proof}

\section{The configuration \texorpdfstring{$\{0,1,3,4,5,7\}\cdot\pi/4$}{A}}

From now on the unit is $\pi/4$, so $\T=\mathbb{R}/8\mathbb{Z}$, and
\[A=\{0,1,3,4,5,7\}.\]
Let $f_k:=\arc{k-1}{k}$, $k=1,\dots,8$, be the faces of $\{0,\dots,7\}$. In the drawing
convention where direction $0$ points up, $\{0,4\}$ is the vertical and $\{2,6\}$ the horizontal
line. The three lines of $A$ are $\{0,4\}$, $\{1,5\}$ and $\{3,7\}$, and the straight subsets of
$A$ are exactly the unions of these. Consider
\begin{align*}
H_1&=\bigl(\{0,4\},\ \{1,7\},\ \{3,5\}\bigr),\\
H_2&=\bigl(\{1,5\},\ \{3,7\},\ \{0,2,4,6\},\ \{2,6\}\bigr),\\
H_3&=\bigl(\{1,5\},\ \{3,7\},\ \{0,4\}\bigr).
\end{align*}

\begin{proposition}\label{prop:base}\leavevmode
\begin{enumerate}[label=(\alph*)]
\item $H_1,H_2,H_3$ are histories of $A$.
\item $\Surv(H_1)=\{\{0,4\},\{1,7\},\{3,5\}\}$: one straight block and two non-straight blocks,
which are mirror images of each other across the line $\{0,4\}$ (since $7=-1$, $5=-3$).
$\Surv(H_2)=\Surv(H_3)=\Lambda:=\{\{0,4\},\{1,5\},\{3,7\}\}$: three straight blocks.
\item $\cost(H_1)=\cost(H_3)=0$ and $\cost(H_2)=2$. In particular $\Lambda$ is the surviving
partition of a history of cost $0$, namely $H_3$.
\item Every step of $H_1$ and $H_2$ satisfies (C). Step $3$ of $H_3$ violates (C).
\end{enumerate}
\end{proposition}
\begin{proof}
We record the normalised folding maps (Definition~\ref{def:flat}) face by face.

\emph{$H_1$.} $\Pi_\emptyset(q(0))=\{0,4\}=D_1$. For $S_1=\{0,4\}$:
$\Phi=x$ on $\arc04$ and $-x$ on $\arc48$, so $\Pi_{S_1}(q(1))=\{1,7\}=D_2$.
For $S_2=\{0,1,4,7\}$, $\Phi$ is $x,\ 2-x,\ x+2,\ -x$ on $\arc01,\arc14,\arc47,\arc78$
respectively, so
$\Phi(3)=\Phi(5)=7$ and $\Pi_{S_2}(q(7))=\{3,5\}=D_3$.

\emph{$H_2,H_3$.} $\Pi_\emptyset(q(1))=\{1,5\}=D_1$. For $S_1=\{1,5\}$: $\Phi=x$ on $\arc15$
and $2-x$ on $\arc59$, so $\Pi_{S_1}(q(3))=\{3,7\}=D_2$. For $S_2=\{1,3,5,7\}$:
$\Phi=x,\ 6-x,\ x-4,\ 2-x$ on $\arc13,\arc35,\arc57,\arc79$; rays map to $1,3,1,3$, and
$\Phi(0)=\Phi(2)=\Phi(4)=\Phi(6)=2$, so $\Pi_{S_2}(q(2))=\{0,2,4,6\}$. Hence $D_3$ is a step for
$H_2$, and so is $\{0,4\}$ for $H_3$. For $S_3=\{0,\dots,7\}$ in $H_2$: $\Phi(k)=0$ for even $k$ and $\Phi(k)=1$ for odd $k$,
with all faces mapping onto $\arc01$, so $\Pi_{S_3}(q(0))=\{0,2,4,6\}\supseteq D_4$.

(b) and (c) are read off directly. For (d): in $H_1$ and in steps $1$--$3$ of $H_2$ we have
$D=\Pi_S(\ell)$, so no face avoiding $D$ contains a preimage of $\ell$. In step $4$ of $H_2$,
$\Pi_{S_3}(q(0))$ consists of rays, so no face contains one in its interior. In step $3$ of
$H_3$, the faces $\arc13$ and $\arc57$ of $S_2$ avoid $\{0,4\}$ and contain $2$ and $6$, with
$\Phi(2)=\Phi(6)=2\in q^{-1}(q(2))$; the labeling switches at $0$ and $4$, so these faces
receive opposite labels.
\end{proof}

\begin{theorem}\label{thm:main}
$H_2$ is not rigidly realizable.
\end{theorem}
\begin{proof}
Steps $1$--$3$ of $H_2$ satisfy the hypothesis of Lemma~\ref{lem:simple}, so each admits exactly
four candidate realizations (two labelings times two choices of $\up$), each determined by
(i)--(iii). Running all $4^3=64$ combinations from the initial layered state (Appendix~\ref{app})
yields, after step $3$, eight faces $f_1,\dots,f_8$ that all have the same image arc $W$
(up to the global isometry $g$), of length $1$ with an endpoint in $q^{-1}(\ell_4)$. Their order $<_W$, bottom to top, is one of
the following eight; the second column is the reversal of the first.
\[
\begin{array}{ll}
f_2f_1f_6f_5\,f_4f_7f_8f_3 & f_3f_8f_7f_4\,f_5f_6f_1f_2\\
f_4f_7f_8f_3\,f_2f_1f_6f_5 & f_5f_6f_1f_2\,f_3f_8f_7f_4\\
f_6f_5f_2f_1\,f_8f_3f_4f_7 & f_7f_4f_3f_8\,f_1f_2f_5f_6\\
f_8f_3f_4f_7\,f_6f_5f_2f_1 & f_1f_2f_5f_6\,f_7f_4f_3f_8
\end{array}
\]
In step $4$ ($D_4=\{2,6\}$) the two labelings move $M=\{f_3,f_4,f_5,f_6\}$ or its complement.
Since $W$ lies in a single component of $\T\setminus q^{-1}(\ell_4)$, condition (F) at any
$u\in W$ says that the moving set is an up-set of $<_W$ (if $W\subseteq\up$) or a down-set
(if $W\subseteq\down$). In each of the eight orders, neither $M$ nor its complement is the set of
the four lowest faces. Hence (F) fails for every choice.
\end{proof}

\begin{remark}
Theorem~\ref{thm:main} does not depend on whether simple folds preserve (O1)--(O3): if they do
not, $H_2$ has no realization at all. It shows that layer data are needed, since $H_2$ passes
(C) by Proposition~\ref{prop:base}(d). We expect $H_1$, a sequence of three simple folds, to be
rigidly realizable, but have not written out the (O1)--(O3) check.
\end{remark}

\begin{conjecture}[Three lines]\label{conj:lines}
$\Lambda\notin\Prig(A)$. That is, no rigidly realizable history of $A$, with any number of
scaffold creases and in any order of steps, has the surviving blocks $\{0,4\}$, $\{1,5\}$,
$\{3,7\}$.
\end{conjecture}

\begin{conjecture}[No straight surviving partition]\label{conj:straight}
No partition in $\Prig(A)$ has all its blocks straight.
\end{conjecture}

\begin{remark}\label{rem:conj}\leavevmode
\begin{enumerate}[label=(\arabic*),leftmargin=*]
\item \emph{Relation.} Conjecture~\ref{conj:straight} implies Conjecture~\ref{conj:lines}. The
partitions of $A$ into straight blocks are the five partitions of the set of its three lines:
$\Lambda$, the three partitions merging two lines, and $\{A\}$.
\item \emph{Order is free.} Both conjectures constrain only $\Surv(H)$, so they cover every order
in which the blocks may be finally created, e.g.\ $\{0,4\}$ first, in the middle, or last.
\item \emph{Symmetry.} The isometries $x\mapsto -x$, $x\mapsto 4-x$ and $x\mapsto x+4$ preserve
$A$ and permute its three lines. The model is equivariant under $\Iso(\T)$ (the normalisation
in Definition~\ref{def:flat} only changes images by a global isometry), so both conjectures are
invariant under these symmetries.
\item \emph{Evidence.} $H_2$ and $H_3$, the two histories with $\Surv=\Lambda$ exhibited here,
are not rigidly realizable: $H_2$ by Theorem~\ref{thm:main}, and $H_3$ because its third step
violates (C), hence (F) by Lemma~\ref{lem:FC}. By contrast $H_1$, which we expect to be rigidly
realizable, has only one straight block.
\item \emph{Status.} Histories with a given surviving partition are unbounded in length, so
neither conjecture is a finite check. A counterexample, if one exists, can be found by
enumerating combinatorial history types.
\item \emph{Rigidity is essential.} Suppose instead that creases persist once made and fold angles
vary continuously, so that unfolding a crease to angle $\pi$ does not erase it. Then $\Lambda$
is realized trivially: fold $\{1,5\}$, fold $\{3,7\}$, unfold both to the flat sheet, fold
$\{0,4\}$, and refold the existing creases as $\{1,7\}$ and $\{3,5\}$, as in $H_1$. Each crease
of $A$ is created exactly once, in the blocks $\{1,5\}$, $\{3,7\}$, $\{0,4\}$. More generally,
connectedness results for single-vertex configuration spaces (Streinu--Whiteley,
\emph{Single-vertex origami and spherical expansive motions}, JCDCG 2004, LNCS 3742; exact
hypotheses not re-checked) suggest that any folded state is then reachable.
\end{enumerate}
\end{remark}

\appendix
\section{Verification script for Theorem~\ref{thm:main}}\label{app}
Unit arcs $\arc{a}{a+1}$ of the image carry stacks of faces (bottom to top). A fold through the
image point $c$ moves the arcs on one side, reflects them via $a\mapsto 2c-a-1$, reverses each
moved stack, and places it on top (over) or at the bottom (under) of the stack at the landing
arc. The fold point of each step is the current image of a crease point of that step.
In \texttt{choices}, \texttt{side} selects the labeling (which side moves) and \texttt{over}
selects $\up$: the moving side is $\up$ iff \texttt{over}, matching rule (iii).
{\footnotesize
\begin{verbatim}
import itertools
def run(choices):
    g = {k: (1, 0) for k in range(1, 9)}        # image map of f_k: x -> e*x+t (mod 8)
    stack = {a: [a+1] for a in range(8)}        # unit arc (a,a+1) -> faces, bottom->top
    img = lambda k, x: (g[k][0]*x + g[k][1]) % 8
    for (pt, face), (side, over) in zip([(1,1), (3,3), (0,1)], choices):
        c = img(face, pt)                       # fold line through image of a crease point
        plus = lambda a: ((a - c) % 8) < 4
        mov = {a for a in range(8) if plus(a) == (side == 0)}
        new = {a: list(stack[a]) for a in range(8) if a not in mov}
        for a in mov:
            b = (2*c - a - 1) % 8
            r = list(reversed(stack[a])); new.setdefault(b, [])
            new[b] = new[b] + r if over else r + new[b]
        for a in mov:
            for k in stack[a]:
                e, t = g[k]; g[k] = (-e, (2*c - t) % 8)
        stack = {a: new.get(a, []) for a in range(8)}
    return stack
M = {3, 4, 5, 6}
res = set()
for ch in itertools.product(itertools.product((0,1),(0,1)), repeat=3):
    st = run(ch); arcs = [a for a in st if st[a]]
    assert len(arcs) == 1 and len(st[arcs[0]]) == 8   # all faces share one image arc
    res.add(tuple(st[arcs[0]]))
for s in sorted(res):                                # 8 orders, none admits a free block
    print(s, set(s[:4]) in (M, set(range(1, 9)) - M))
\end{verbatim}
}
The output consists of the eight orders listed in the proof of Theorem~\ref{thm:main}, each with
the value \texttt{False}.

\end{document}
